We first record the horseshoe class and the conditional genericity statement,
then discuss why the hypothesis of (ii) is delicate.

\begin{definition}\label{def:O-recall}
$\mathcal O\subseteq\Cs$ is the set of maps carrying a \emph{strict two-branch
horseshoe}: disjoint nonempty closed intervals $A_1<A_2$ and $\theta>0$ with
$g(A_i)\supseteq[\min A_1-\theta,\,\max A_2+\theta]\subseteq(0,1)$ for $i=1,2$.
As in \cite{companion}, $\mathcal O$ is $\rsup$-open (an interval image
containing an interval with margin $\theta$ persists with margin $\theta/2$
under $\rsup$-perturbations below $\theta/2$), every $g\in\mathcal O$ has
topological entropy at least $\log2$, and every map in $\mathcal O$ is DC1 by
Schweizer--Sm\'\i tal \cite{SS}.
\end{definition}

\begin{proof}[Proof of Theorem \ref{thm:strata}(ii)]
By (i) we may work with $\rsup$ and its system metric, and $\mathcal F^\mu$ is a
complete metric space.  The limit-map function
$\lm(f_{1,\infty})=\lim_nf_n$ is $1$-Lipschitz from $\mathcal F^\mu$ to
$(C^\mu_s,\rsup)$, so
$U:=\lm^{-1}(\mathcal O\cap C^\mu_s)\cap\mathcal F^\mu$ is relatively open.
It is dense: given $(f_n)\in\mathcal F^\mu$ with limit $f$ and
$\varepsilon>0$, the hypothesis of (ii) provides
$h\in\mathcal O\cap C^\mu_s$ with $\rsup(h,f)<\varepsilon/2$; pick $N$ with
$\rsup(f_n,f)<\varepsilon/2$ for $n>N$ and pass to the system
$(f_1,\dots,f_N,h,h,\dots)\in\mathcal F^\mu$, which is $\varepsilon$-close and
has limit $h$.  Every system in $U$ is uniformly convergent (by (i)) to a limit
map carrying a strict horseshoe, hence of entropy at least $\log2$, hence DC1 by
\cite{SS}; by \v Stef\'ankov\'a's inheritance theorem \cite{Stefankova} the
system itself is DC1.  Thus the DC1 systems contain the relatively open dense
set $U$ and are residual in $\mathcal F^\mu$.
\end{proof}

\subsection*{The horseshoe-density hypothesis: two obstructions and an open
problem}

For the supremum stratum --- no modulus constraint --- the density of
$\mathcal O$ in $\Cs$ was proved in \cite{companion} by implanting a steep
zigzag near a fixed point.  Inside a stratum $C^\mu_s$ the same strategy meets
two quantitative obstructions, which we record because they delimit any future
proof.

\emph{Saturation.}  The implanted map $v$ differs from the base map $u$ by a
perturbation $\zeta$ supported near the implantation site.  For a pair $(x,y)$
with $x$ in the site and $y$ far from it,
$|v(x)-v(y)|\le|u(x)-u(y)|+\|\zeta\|_\infty$, and if $u$ \emph{saturates} its
modulus along such pairs ($|u(x)-u(y)|=\mu(|x-y|)$, which mollification does not
destroy), every nonconstant perturbation overruns the budget.  Any density proof
must therefore first create slack at all scales --- for instance by the
rescaling $u\mapsto(1-\gamma)u+\gamma/2$ --- and such rescalings destroy
surjectivity.

\emph{The repair budget.}  Restoring surjectivity after a rescaling by
$1-\gamma$ requires moving some value back to $0$ (and to $1$).  Within
$C^\mu$ a correction of height $\gamma/2$ built from the profile
$c\,\mu(|x-x_0|)$ needs $c\,\mu(\operatorname{radius})\ge\gamma/2$; the modulus
budget available for corrections is of order $\gamma\mu$, so the correction
radius must satisfy $\mu(\mathrm{radius})\gtrsim\tfrac12$, i.e.\ the radius is
of order $\mu^{-1}(\tfrac12)$ --- \emph{not} small.  Large-support corrections
interact with the implanted horseshoe and with each other, and for moduli with
$\mu(1)$ close to $1$ they may not fit in $[0,1]$ at all.  The two requirements
--- slack everywhere, surjectivity at the endpoints --- thus compete for the
same budget.

These obstructions are quantitative, not absolute: they disappear, for example,
on strata where every map admits a point of strict sub-saturation at the
relevant scales.  We do not know the answer in general and consider the
following the main question left open by this paper.

\begin{problem}\label{prob:strata}
For which moduli $\mu$ is the strict-horseshoe class $\mathcal O\cap C^\mu_s$
dense in $(C^\mu_s,\rsup)$ --- equivalently (by
Theorem~\ref{thm:strata}(ii)), for which strata is DC1 generic in
$\mathcal F^\mu$?  By Theorem~\ref{thm:strata}(iii) the answer is negative for
$\mu(t)\le t$.  Is it positive for every non-Lipschitz concave modulus?  For
Lipschitz strata $\mu(t)=Lt$, is it positive for every $L>1$, with horseshoes
produced at the level of an iterate (where the Lipschitz budget compounds to
$L^n$), or supplied by a crookedness dichotomy (maps saturating the budget at
every scale should already carry horseshoes)?
\end{problem}
